% Small print among normal text: \small to \tiny, a footnote, a table, a caption, a margin note.
% pdflatex smallprint.tex makes smallprint.pdf.
\documentclass[11pt]{article}
\usepackage[utf8]{inputenc}
\usepackage[T1]{fontenc}
\usepackage[ngerman]{babel}
\usepackage{times}
\begin{document}
Das Velo steht auf dem Trottoir, und die Offerte liegt im Estrich.\footnote{Das Velo im Spital wurde parkiert; die Glace war gross.}

{\small Kleine Schrift: Das Velo steht vor dem Spital.}

{\footnotesize Fussnotengrösse: Wir parkieren das Velo beim Trottoir.}

{\scriptsize Sehr klein: Die Offerte für das Velo ist gross.}

{\tiny Winzig: Das Trottoir beim Spital ist gross.}

\begin{table}[h]\centering
\begin{tabular}{ll}
Velo & Trottoir \\
Offerte & Spital \\
{\scriptsize Velostation} & {\tiny Glace} \\
\end{tabular}
\caption{Velos und Trottoirs im Kanton}
\end{table}

\begin{center}{\scriptsize Velo}\quad{\tiny Velo}\quad Velo\textsuperscript{2} und H$_{2}$O beim Spital\end{center}

Normal: Der Entscheid über die Strasse ist gross.\marginpar{\tiny Velo am Rand}
\end{document}
